\section*{Chapter \ref{ch:g11:fossil-fuel-geology} --- The Geology of Coal, Oil and Gas}

\begin{solution}{exo:g11:fossil-fuel-geology:1}
Dead plankton settles on a sea floor poor in oxygen; the organic matter
becomes \omterm{def:g11:fossil-fuel-geology:source-rock}{kerogen} in a mud; burial under younger sediments; \omterm{def:g11:fossil-fuel-geology:oil-window}{maturation} in the
\omterm{def:g11:fossil-fuel-geology:oil-window}{oil window}; migration; oil gathers under a \omterm{def:g11:fossil-fuel-geology:reservoir}{seal}.
\end{solution}

\begin{solution}{exo:g11:fossil-fuel-geology:2}
(1) \omterm{def:g11:fossil-fuel-geology:source-rock}{Source rock}. (2) \omterm{def:g11:fossil-fuel-geology:reservoir}{Reservoir rock}. (3) \omterm{def:g11:fossil-fuel-geology:reservoir}{Seal}. (4) \omterm{def:g11:fossil-fuel-geology:reservoir}{Seal}. (5) \omterm{def:g11:fossil-fuel-geology:reservoir}{Reservoir rock}.
\end{solution}

\begin{solution}{exo:g11:fossil-fuel-geology:3}
Top: $(60 - 10)/25 = \qty{2.0}{km}$; base: $(120 - 10)/25 = \qty{4.4}{km}$.
The window is $60/25 = \qty{2.4}{km}$ thick.
\end{solution}

\begin{solution}{exo:g11:fossil-fuel-geology:4}
(1) Peat, lignite, bituminous \omterm{def:g11:fossil-fuel-geology:coal}{coal}, anthracite. (2) Anthracite, 86 to
\qty{97}{\percent} carbon. (3) Anthracite, beyond about \qty{5}{km} and
\qty{150}{\degreeCelsius}.
\end{solution}

\begin{solution}{exo:g11:fossil-fuel-geology:5}
(1) The fluids sort themselves by density: gas is the lightest, then oil,
then water. (2) Oil and gas rise because they are lighter than water; only
a rock they cannot pass through stops them. Mudstone (or shale) and salt.
\end{solution}

\begin{solution}{exo:g11:fossil-fuel-geology:6}
(1) A: $(60 - 15)/20 = \qty{2.25}{km}$ and $(120 - 15)/20 = \qty{5.25}{km}$.
B: $45/40 \approx \qty{1.13}{km}$ and $105/40 \approx \qty{2.63}{km}$.
(2) At \qty{2}{km}: in A, $15 + 40 = \qty{55}{\degreeCelsius}$, not yet
mature, no oil; in B, $15 + 80 = \qty{95}{\degreeCelsius}$, in the \omterm{def:g11:fossil-fuel-geology:oil-window}{oil
window}, it makes oil. (3) In basin A: oil forms there only below
\qty{2.25}{km}.
\end{solution}

\begin{solution}{exo:g11:fossil-fuel-geology:7}
(1) $11 \times 3 = \qty{33}{m}$ of peat. (2) $\qty{33000}{mm}$ at
\qty{1}{mm} a year: about \num{33000} years. (3) The peat stays only where
it is under stagnant water, out of reach of oxygen; as it piles up, the
ground must sink at the same pace to keep its surface waterlogged.
\end{solution}

\begin{solution}{exo:g11:fossil-fuel-geology:8}
(1) $20 + 30 \times 1.2 = \qty{56}{\degreeCelsius}$;
$20 + 30 \times 2.3 = \qty{89}{\degreeCelsius}$;
$20 + 30 \times 4.0 = \qty{140}{\degreeCelsius}$.
(2) 120 million years ago, immature; 95 million years ago, in the \omterm{def:g11:fossil-fuel-geology:oil-window}{oil
window}; 20 million years ago, in the \omterm{def:g11:fossil-fuel-geology:oil-window}{gas window}.
\end{solution}

\begin{solution}{exo:g11:fossil-fuel-geology:9}
(1) $2580 - 2400 = \qty{180}{m}$. (2) Gas fills
$2450 - 2400 = \qty{50}{m}$, oil $2580 - 2450 = \qty{130}{m}$. (3) They spill
out at the spill point, below the deepest closed contour, and migrate on
upwards, to another \omterm{def:g11:fossil-fuel-geology:trap}{trap} or to the surface.
\end{solution}

\begin{solution}{exo:g11:fossil-fuel-geology:10}
(1) Top of the window at \qty{3}{km}: $(60 - 20)/3 \approx
\qty{13}{\degreeCelsius/km}$; base: $(120 - 20)/3 \approx
\qty{33}{\degreeCelsius/km}$. (2) At \qty{3}{km} the temperature is
$20 + 75 = \qty{95}{\degreeCelsius}$ to $20 + 102 = \qty{122}{\degreeCelsius}$:
in the \omterm{def:g11:fossil-fuel-geology:oil-window}{oil window}, or at its base for the steepest gradients. The rule is
reasonable: in most basins a \omterm{def:g11:fossil-fuel-geology:source-rock}{source rock} under \qty{3}{km} of rocks has
started to make oil.
\end{solution}

\begin{solution}{exo:g11:fossil-fuel-geology:11}
(1) \qty{30}{\percent}: lignite; \qty{70}{\percent}: bituminous \omterm{def:g11:fossil-fuel-geology:coal}{coal};
\qty{92}{\percent}: anthracite. (2) The anthracite, beyond about
\qty{5}{km} and \qty{150}{\degreeCelsius}. (3) Black and glassy.
\end{solution}

\begin{solution}{exo:g11:fossil-fuel-geology:12}
(1) $1100/25 = \qty{44}{m}$ per million years. (2) $0.044 \times 30
\approx \qty{1.3}{\degreeCelsius}$ per million years. (3) The window spans
\qty{60}{\degreeCelsius}: $60/1.32 \approx 45$ million years. (4) Burial
slowed down after 95 million years ago (about 24, then 23 and
\qty{20}{m} per million years): the \omterm{def:g11:fossil-fuel-geology:source-rock}{source rock} warmed more slowly.
\end{solution}

\begin{solution}{exo:g11:fossil-fuel-geology:13}
(1) $(80 - 51)/(117 - 51) = 29/66 \approx 0.44$: about \qty{44}{\percent}.
(2) It migrated upwards before any \omterm{def:g11:fossil-fuel-geology:trap}{trap} existed above: it escaped to the
surface, where it was oxidised, or gathered in some other \omterm{def:g11:fossil-fuel-geology:trap}{trap}. (3) By 40
million years ago the \omterm{def:g11:fossil-fuel-geology:source-rock}{source rock} had left the \omterm{def:g11:fossil-fuel-geology:oil-window}{oil window}: it made only gas,
and the anticline could hold only gas from it.
\end{solution}

\begin{solution}{exo:g11:fossil-fuel-geology:14}
(1) Top: $(60 - 20)/40 = \qty{1.0}{km}$; base: $(120 - 20)/40 =
\qty{2.5}{km}$. (2) \qty{1.0}{km} lies between \qty{0.4}{km} (150 million
years ago) and \qty{1.2}{km} (120): $150 - 30 \times 0.6/0.8 = 127.5$
million years ago. \qty{2.5}{km} lies between \qty{2.3}{km} (95) and
\qty{3.0}{km} (66): $95 - 29 \times 0.2/0.7 \approx 86.7$ million years
ago. (3) No oil: the \omterm{def:g11:fossil-fuel-geology:source-rock}{source rock} left the \omterm{def:g11:fossil-fuel-geology:oil-window}{oil window} before the anticline
formed, and its oil had migrated away. It could hold gas made after 80
million years ago.
\end{solution}

\begin{solution}{exo:g11:fossil-fuel-geology:15}
(1) At least \qty{5}{km}. (2) $15 + 30 \times 5 = \qty{165}{\degreeCelsius}$,
above \qty{150}{\degreeCelsius}: consistent. (3) At least \qty{5}{km} of
rock has been uplifted and eroded. $\qty{5e6}{mm}$ in at most
\num{299} million years gives at least $5 \times 10^6/2.99 \times 10^8
\approx \qty{0.017}{mm}$ a year on average.
\end{solution}

\begin{solution}{pb:g11:fossil-fuel-geology:1}
\textbf{1.} \omterm{def:g11:fossil-fuel-geology:source-rock}{Source rock}: the black shale; reservoir: the porous sandstone;
\omterm{def:g11:fossil-fuel-geology:reservoir}{seal}: the mudstone above it.

\textbf{2.} Plankton that settled on a sea floor poor in oxygen escaped
decay and was buried in the mud.

\textbf{3.} The \omterm{def:g11:fossil-fuel-geology:source-rock}{source rock} in the Jurassic (165 lies between 201.4 and
143.1). The sandstone began in the Jurassic and was mostly laid down in the
Early Cretaceous.

\textbf{4.} The organic matter scattered in a sedimentary rock, insoluble in
organic solvents. It makes oil only once heated to about
\qty{60}{\degreeCelsius}, which takes burial.

\textbf{5.} $400 + 800 + 1100 + 700 + 700 + 300 - 500 = \qty{3500}{m}$.

\textbf{6.} \qty{0.4}{km}, \qty{1.2}{km}, \qty{2.3}{km}, \qty{3.0}{km},
\qty{3.7}{km}, \qty{4.0}{km}; today \qty{3.5}{km}.

\textbf{7.} $400/15 \approx 26.7$; $800/30 \approx 26.7$; $1100/25 = 44$;
$700/29 \approx 24.1$; $700/31 \approx 22.6$; $300/15 = 20$ (\unit{m} per
million years).

\textbf{8.} 20 million years ago, at \qty{4.0}{km}.

\textbf{9.} The basin was uplifted and \qty{500}{m} of rock eroded from its
top.

\textbf{10.} The sandstone's base was laid down 150 million years ago, on
top of \qty{400}{m} of shale and mudstone: $3.5 - 0.4 = \qty{3.1}{km}$.

\textbf{11.} $T = 20 + 30 z$: \qty{32}{\degreeCelsius}, \num{56}, \num{89},
\num{110}, \num{131}, \qty{140}{\degreeCelsius}; today
\qty{125}{\degreeCelsius}.

\textbf{12.} $40/30 \approx \qty{1.33}{km}$ and $100/30 \approx
\qty{3.33}{km}$.

\textbf{13.} Between 120 million years ago (\qty{56}{\degreeCelsius}) and
95 million years ago (\qty{89}{\degreeCelsius}).

\textbf{14.} The temperature rose by 33 degrees in 25 million years; it took
$25 \times 4/33 \approx 3.0$ million years to rise by the 4 degrees from
\num{56} to \qty{60}{\degreeCelsius}: about $120 - 3.0 = 117$ million years
ago.

\textbf{15.} Between 66 million years ago (\qty{110}{\degreeCelsius}) and
35 (\qty{131}{\degreeCelsius}): $66 - 31 \times 10/21 \approx 51.2$ million
years ago.

\textbf{16.} \qty{140}{\degreeCelsius}, 20 million years ago. At
\qty{125}{\degreeCelsius} today it is in the \omterm{def:g11:fossil-fuel-geology:oil-window}{gas window}: it makes mostly
methane.

\textbf{17.} With no \omterm{def:g11:fossil-fuel-geology:trap}{trap} above it, it migrated up and escaped, to the
surface or to other \omterm{def:g11:fossil-fuel-geology:trap}{traps}.

\textbf{18.} $(80 - 51.2)/(117.0 - 51.2) \approx 0.44$.

\textbf{19.} Gas at the top, then oil, then water.

\textbf{20.} The \omterm{def:g11:fossil-fuel-geology:source-rock}{source rock} entered the \omterm{def:g11:fossil-fuel-geology:oil-window}{oil window} about \textbf{117
million years ago}.
\end{solution}
